

To create a projective variety, we can utilize the command \verb'-projective_variety'. 
In order to do so, we first need to define a polynomial ring and an associated projective space. 
We use a symbolic object to define the equation. 
In the following command we create the Fermat curve as a projective variety. 
The difference to the command before is that now the curve is stored as an object of type projective variety. 
In the previous example, the equation of the curve was parsed, but the curve was not remembered as an object.
Here is the command to create the variety for a Fermat cubic in $\PG(2,9)$. 

\sourcecode{Fermat_cubic_F9}


\bigskip



Next, we will consider ideals. As an application, we will 
classify arcs in a projective plane and see which conics we get.
The next command classifies the $(5,2)$-arcs in $\PG(2,11)$:

\sourcecode{arcs_5_2_q11}

It finds exactly two isomorphism types of arcs. 
The representative sets are 
$$
\{0, 1, 2, 3, 37\}, \quad 
\{ 0, 1, 2, 3, 49\}.
$$
They are stored in the file \verb'arcs_5_2_q11_lvl_5'. Let us now create the ideal 
in the quadratic component of the polynomial ring in three variables over $\bbF_{11}$:

\sourcecode{arcs_5_2_q11_ideal}

The ideals are generated by  
$$
\verb'7*x0*x1 + 5*x0*x2 + 10*x1*x2'
$$
and
$$ 
\verb'4*x0*x1 + 8*x0*x2 + 10*x1*x2',
$$
respectively.


\bigskip


Let us consider a smooth cubic surface with 9 lines and 4 Eckardt points.
Suppose we have the set of points and we wish to determine the equation of the object. 
To do so, we first define the object from the given set of points. 

\sourcecode{PTS_OF_SURFACE_ORBIT211_Q3_L9_E4}

Then, we create a ring and compute the ideal:

\sourcecode{surface_9lines_4E_ideal}

We find a two-dimensional ideal. Generators are:
$$
\verb'x0*x0*x1 + 2*x0*x1*x1 + 2*x0*x1*x3'
\quad\mbox{and}\quad
\verb'2*x2*x2*x3 + 2*x2*x3*x3'.
$$
Let us take the sum of the two polynomials and create the cubic surface:

\sourcecode{SURFACE_F_9}


\sourcecode{F_9_q7}



\bigskip

In the next example, 
we wish to explore the relationship between conics and $(5,2)$-arcs. 
We consider the plane $\PG(2,11).$
Instead of classification, we will try random generation this time. 
Since there are 133 points, we create a number of 5-subsets of a set of size 133. 
In this case, we create 20 sets at random:

\sourcecode{random_k_subsets_PG_2_11}

The sets are stored in the file \verb'random_k_subsets_n133_k5_nb20.csv'.
Now, let's compute the line type of these subsets, to see which ones are arcs:

\sourcecode{line_type_in_PG_2_11}

It turns out that the second set is an arc. It is the set $\{3,33,40,83,102\}.$
We create the conic through these 5 points:

\sourcecode{random_arc_5_2_q11_ideal}

The ideal is generated by 
$$
\verb'10*x0*x0 + 3*x0*x1 + 8*x0*x2 + 2*x1*x1 + 10*x2*x2'.
$$
The conic contains the following 12 points:
$$
\{3, 15, 19, 33, 40, 42, 46, 50, 83, 88, 102, 108\}.
$$


\bigskip




The next command creates the Endrass surface over $\bbF_7$. 
The surface is defined as a makefile variable in sparse form.

\sourcecode{ENDRASS_SPARSE}


\sourcecode{Endrass_F7.txt}



\bigskip


Suppose we want to create the monomials of degree 8 in 4 variables. 
We use an diophantine system to do so. 
The following command creates the system and solves it. 
After that, it applies the unix sort command to sort the monomials:

\sourcecode{octic_prepare}

There are 165 monomials. 
They are listed in the file \verb'octic_monomials_sorted.txt'.


\bigskip


